% Copyright (c) 2026 Eugene Joseph M. Ragasa.
% SPDX-License-Identifier: Apache-2.0
% Author: Eugene Joseph M. Ragasa
% Formal proof completed: 2026-10-19
\documentclass[11pt]{article}
\usepackage[T1]{fontenc}\usepackage{amsmath,amssymb}\usepackage[margin=1in]{geometry}\usepackage[hidelinks]{hyperref}
\title{Proof Dossier: Ragasa2026A Canonical Spectral Power}\author{Eugene Joseph M. Ragasa}\date{Formal proof completed: 2026-10-19}
\begin{document}\maketitle
\section{Identity and status}
ID \texttt{BG-RAGASA2026A-THM-019}; status \texttt{proved}. Lean declaration:
\par\noindent {\tiny\path{ProjectKoios.BandGauge.Papers.Ragasa2026A.canonicalSpectralSplittingField_power_eq_parent}}.
\section{Exact statement}
For the canonical Hermitian spectral decomposition,
\[\sum_i\frac{\alpha_i}{2E_{\mathrm{ref}}^2}\sum_n d_{i,n}^2
 =\sum_i\frac{\alpha_i}{2E_{\mathrm{ref}}^2}\lVert D(q_i)\rVert_F^2.\]
\section{Source, proof, and boundary}
This \texttt{formalizes} the common-power identification used in manuscript lines
603--610 at owner commit
\begin{center}\small\texttt{99a46cb59d0c6f27ada7d709f194996ef9619d9f}.\end{center}
It specializes unitary Frobenius invariance to the canonical spectral-theorem witness.
No numerical eigenvalues are produced. Kernel axioms: \texttt{propext},
\texttt{Classical.choice}, \texttt{Quot.sound}; no admissions. AI-assisted and
kernel-checked; no independent expert review recorded.
\end{document}
